% !TeX root = ../main.tex
\chapter{Experimental Design}
\label{ch:experimental-design}
\thesisplaceholder{Define the evaluation protocol, selection procedure, uncertainty analysis, and robustness checks.}

\section{Data Splits and Leakage Prevention}
% 根據問題選擇病人、donor、批次、院所或時間切分，保留不可接觸的測試資料。
% 檢查同一個體或來源是否跨 split；外部驗證須有獨立來源依據。

\section{Model Selection and Fair Comparisons}
% 事先說明調參範圍、驗證資料、計算預算、選模準則與測試集使用次數。
% 所有方法使用可比較的資料與相同終點；記錄任何預訓練重疊風險。

\section{Metrics and Statistical Analysis}
% 定義主要／次要指標、估計目標、閾值、統計單位與不確定性估計。
% 依依賴結構處理重複觀測；說明多重比較、缺失值與樣本量的限制。

\section{Robustness, Subgroups, and Ablations}
% 列出與實際用途相關的亞組、分布改變、消融及敏感度分析。
% 預先規劃與事後探索分開；解釋小亞組的不確定性，勿過度解讀。
